\documentclass[11pt]{article}
\usepackage[a4paper,margin=22mm]{geometry}
\usepackage[T1]{fontenc}
\usepackage{lmodern}
\usepackage{amsmath}
\usepackage{microtype}
\usepackage{graphicx}
\usepackage{xcolor}
\usepackage{float}
\usepackage{fancyhdr}
\usepackage{booktabs}
\usepackage[hidelinks]{hyperref}

\definecolor{Accent}{HTML}{21918C}
\definecolor{Muted}{HTML}{5B6570}
\definecolor{Warn}{HTML}{B03A2E}

\newcommand{\Faq}{\textit{F.~aquilonia}}
\newcommand{\Fpo}{\textit{F.~polyctena}}

\pagestyle{fancy}\fancyhf{}
\lhead{\small\color{Muted}Ancestry, climate \& mitotype -- supplementary analysis (corrected)}
\rhead{\small\color{Muted}Poikela et al., rho05 pipeline}
\cfoot{\small\thepage}
\renewcommand{\headrulewidth}{0.4pt}
\setlength{\parskip}{0.5em}\setlength{\parindent}{0pt}

\title{\vspace{-3em}\textbf{\color{Accent}Genome-wide ancestry, climate, and mitotype}\\[0.2em]
\large A population-level follow-up to the BayPass climate-GEA scan \\[0.3em]
\normalsize\color{Warn}Corrected version -- see Erratum}
\author{}
\date{\small Generated \today}

\begin{document}
\maketitle
\vspace{-2em}

\section*{\color{Warn}Erratum}
An external audit of the first version of this document identified two implementation
bugs, both now fixed. All numbers, tables and figures below are corrected;
Section~\ref{sec:limitations} also incorporates the audit's remaining recommendations.
\begin{enumerate}
  \item \textbf{Block permutation silently dropped one population per shared-origin
    pair.} The reduction from 19 populations to 17 independent lineage units was
    implemented as ``take the first population's value, broadcast it to both members of
    the pair.'' This was harmless for L{\aa}ngholmenW/L{\aa}ngholmenR (identical climate
    values) but wrong for Bunkkeri/Grundsund, whose PC1 values differ ($-3.409$ vs.\
    $-3.899$) -- Grundsund's actual value was discarded, not averaged in. \textbf{Fix:}
    ancestry, climate, and mitotype are now genuinely \emph{averaged} within each
    shared-origin pair to build a real 17-observation dataset, and every correlation,
    null, and permutation below operates on that dataset directly.
  \item \textbf{The $\Omega$-structured null for the partial correlation replaced the
    wrong variable.} $\Omega$ is a nuclear allele-frequency covariance matrix, so it is
    an appropriate null generator only for the one tested quantity that is itself a
    direct function of nuclear allele frequencies -- ancestry. The original
    implementation drew random vectors standing in for PC1 (holding ancestry and
    mitotype fixed) instead of standing in for ancestry (holding PC1 and mitotype
    fixed). \textbf{Fix:} throughout this document, $\Omega$-null draws always replace
    ancestry. For comparisons where \emph{neither} variable is ancestry (mitotype vs.\
    PC1/PC2/bio6/bio11), the $\Omega$-null is not well-motivated and is omitted; those
    rows report the block-permutation result only.
\end{enumerate}
$\Omega$ itself was reduced to match the 17-unit dataset via $M\Omega M^\top$, where $M$
is the same population-to-lineage-unit averaging operator applied to the phenotypes --
propagating the actually-estimated relatedness structure through the reduction, rather
than drawing null vectors from an arbitrary 17-dimensional matrix. Corrected results were
cross-checked against an independent re-implementation: all $r$ values matched exactly,
and $p$-values agreed closely (small differences are expected Monte Carlo noise between
independent permutation runs).

\section*{\color{Accent}Motivation}

The Stage-1-direct BayPass climate-association scan (PC1, PC2, and the winter-temperature
variables bio6/bio11) produced raw candidate counts under a descriptive BF(dB)$\geq$15
threshold, but none of the climate axes are expected to survive full structured-null
calibration cleanly, because PC1 and PC2 are both partially confounded with genome-wide
ancestry (established in \texttt{moduleB\_ancestry\_confound.R}: cor(PC1,ancestry)$=-0.23$,
cor(PC2,ancestry)$=+0.57$ across all 20 populations). Rather than treating that confound
purely as a caveat that explains away the locus-level scan, this document asks whether the
population-level ancestry--climate link is itself real and worth reporting on its own
terms -- and traces it down to what actually drives it.

\textbf{Scope note.} The BayPass scan referenced throughout this document is the
\emph{Stage-1 large-cluster scan}: the covariance basis ($\Omega$) is estimated from all
698{,}251 Stage-1 LD clusters (including singletons), but the tested/calibrated unit
universe is the 18{,}361 clusters with $\geq$5 SNPs (best-SNP represented) -- 2.6\% of all
Stage-1 clusters. Isolated SNPs and smaller clusters were not tested by this scan; the
locus-level conclusions below (Section~5, floor-survivor counts) are scoped accordingly
and should not be read as a statement about the complete 661{,}386-unit LD-reduced genome
(that complete set is used elsewhere in this pipeline, e.g.\ the Fst-vs-DI analysis, but
not for this climate calibration).

All analyses here operate on the population level (\emph{not} individual genotypes), on the
\texttt{aland\_excluded} set of hybrid populations that the BayPass scan itself uses
(\AA land excluded because its individuals carry mixed mitotypes and it is a known
high-leverage outlier in several of the ORIGINAL PC1/PC2 diagnostics), further reduced to
\textbf{17 independent lineage units} as described in the Erratum.

\subsection*{\color{Muted}Two non-independence-aware significance tests}
A plain Pearson correlation across populations overstates the effective sample size: two
population pairs have a documented shared origin from the phylogenetic network analysis
(L{\aa}ngholmenW+L{\aa}ngholmenR; Bunkkeri+Grundsund; Nouhaud et al.\ 2022), and all
populations share allele-frequency covariance captured by the fixed population covariance
matrix $\Omega$ already estimated for the BayPass scan (from all 698{,}251 Stage-1 cluster
representatives). Two complementary null designs are used throughout, both now operating on
the genuine 17-lineage-unit dataset:
\begin{itemize}
  \item \textbf{$\Omega$-structured null.} 10{,}000 random vectors are drawn from the
    17-unit-reduced $\Omega$'s eigendecomposition, standing in for ancestry (see Erratum).
    Empirical $p = (1+k)/(n_{\mathrm{sim}}+1)$, $k$ = number of null draws with
    $|r_{\mathrm{null}}| \geq |r_{\mathrm{obs}}|$.
  \item \textbf{Shared-origin block permutation.} 10{,}000 permutations shuffle the tested
    variable across the 17 independent units directly, recomputing the correlation each
    time.
\end{itemize}
Benjamini-Hochberg correction is applied within each family of tests (the 4 climate
variables; the 5 mitotype comparisons), using the block-permutation $p$-values as the
primary statistic, since the $\Omega$-null is not defined for every row.

\clearpage
\section*{\color{Accent}1. Ancestry vs.\ climate axes}

Population ancestry (mean oriented \Faq\ allele frequency at diagnostic markers,
DI$>-25$, $|\Delta p|\geq0.5$) is compared against the two original climate axes (PC1, PC2)
and the two new winter-temperature variables suggested for the GEA scan
(bio6 = min.\ temperature of the coldest month; bio11 = mean temperature of the coldest
quarter).

\begin{table}[H]\centering\small
\begin{tabular}{lrrrr}
\toprule
Variable & $r$ (17 units) & $\Omega$-null $p$ & Block-perm.\ $p$ & Block-perm.\ $p_{\mathrm{BH}}$ \\
\midrule
PC1   & $-0.481$ & $0.044$ & $0.051$ & $0.164$ \\
PC2   & $+0.381$ & $0.110$ & $0.123$ & $0.164$ \\
bio6  & $+0.404$ & $0.091$ & $0.110$ & $0.164$ \\
bio11 & $+0.287$ & $0.210$ & $0.267$ & $0.267$ \\
\bottomrule
\end{tabular}
\caption{Ancestry vs.\ climate, 17 independent lineage units. $p_{\mathrm{BH}}$ =
Benjamini-Hochberg corrected across the 4 tests.}
\end{table}

\textbf{Finding.} PC1 has the strongest ancestry association among the four variables and
is the only one to fall below $p=0.05$ under the block-permutation test ($p=0.051$, i.e.\
right at the boundary) -- but the two null designs disagree (\smash{$\Omega$-null $p=0.044$}
vs.\ block-permutation $p=0.051$), and \textbf{none of the four variables survive
Benjamini-Hochberg correction} (minimum $p_{\mathrm{BH}}=0.164$). The accurate summary is
that PC1's ancestry association is \emph{suggestive but not formally significant}, not that
it ``clears'' a defensible statistical test. bio6 is the next-closest candidate but is
weaker on both nulls than PC1.

\begin{figure}[H]\centering
  \includegraphics[width=\linewidth]{../Figures/moduleB_ancestry_vs_climate_17units_CORRECTED.png}
  \caption{Ancestry vs.\ PC1/PC2/bio6/bio11, 17 independent lineage units, with corrected
  statistics annotated per panel.}
\end{figure}

\clearpage
\section*{\color{Accent}2. What drives the PC1--ancestry signal?}

Three complementary diagnostics identify which population(s) are responsible for the
PC1--ancestry correlation: leave-one-population-out, leave-one-independent-unit-out, and
standard regression influence statistics (Cook's distance, standardized residuals) from
$\mathrm{lm}(\mathrm{ancestry}\sim\mathrm{PC1})$. \emph{This diagnostic predates the
17-unit correction and is reported on the original 19-population scale} (observed
$r=-0.502$) -- Heinam\"aki and Jarvenpaa are not part of either shared-origin pair, so the
bug did not affect their values, and the qualitative finding is unchanged at 17-unit
resolution (raw $r=-0.481$; see Section~1).

\begin{table}[H]\centering\small
\begin{tabular}{lrrrr}
\toprule
Population & Leave-one-out $\Delta r$ & $r$ without it & Cook's $D$ & Std.\ residual \\
\midrule
Heinam\"aki & $-0.211$ & $-0.712$ & $\mathbf{0.618}$ & $+2.75$ \\
Jarvenpaa   & $+0.030$ & $-0.471$ & $0.288$ & $-2.30$ \\
\emph{all others} & $|{\leq}0.06|$ & -- & $\leq0.07$ & $|{\leq}1.2|$ \\
\bottomrule
\end{tabular}
\caption{Regression influence diagnostics for PC1 vs.\ ancestry, 19-population scale
(observed $r=-0.502$). Rule-of-thumb cutoffs: Cook's $D > 4/n = 0.211$.}
\end{table}

\textbf{Finding.} Heinam\"aki dominates every diagnostic by a wide margin, with Jarvenpaa a
distant second; every other population's individual influence is negligible. Critically,
\textbf{Heinam\"aki dampens rather than drives} the correlation -- it combines high
ancestry ($0.651$) with an extreme PC1 value ($4.29$) in a way that sits \emph{off} the
negative trend, so excluding it makes the correlation \emph{stronger}
($r:-0.502\rightarrow-0.712$). Whatever one makes of PC1's borderline significance (Section
1), it is not an artifact of one population manufacturing a spurious signal.

\begin{figure}[H]\centering
  \includegraphics[width=0.75\linewidth]{../Figures/moduleB_ancestry_PC1_leverage.png}
  \caption{PC1 vs.\ ancestry (19-population scale), point size $=$ Cook's distance.
  Heinam\"aki (largest point) is the dominant influence but pulls \emph{against} the fitted
  trend.}
\end{figure}

\subsection*{\color{Muted}Is Heinam\"aki itself anomalous?}
Three checks on Heinam\"aki's own data found nothing unusual: mitotype is uniform
(10/10 \Faq-like, not mixed like \AA land); it was sampled from 2 nests (159, 160), typical
of the panel; and its hybrid index ($0.357\pm0.017$) and heterozygosity ($0.314\pm0.038$)
sit comfortably among the other \Faq-skewed mainland populations. Admixture timing (DATES)
could not be checked here (not part of this repository).

\clearpage
\section*{\color{Accent}3. Mitotype vs.\ ancestry and climate}

Mitotype (\Faq-like vs.\ \Fpo-like maternal lineage) is uniform within both shared-origin
pairs, so it survives the 17-unit reduction unambiguously (6 \Faq-like, 11 \Fpo-like
units). Tested as a population-level binary variable (point-biserial correlation) against
the same five quantities.

\begin{table}[H]\centering\small
\begin{tabular}{lrrrr}
\toprule
Variable & $r$ with mitotype & $\Omega$-null $p$ & Block-perm.\ $p$ & Block-perm.\ $p_{\mathrm{BH}}$ \\
\midrule
ancestry & $+0.490$ & $0.037$ & $0.049$ & $0.243$ \\
PC1      & $-0.299$ & -- & $0.255$ & $0.424$ \\
PC2      & $+0.197$ & -- & $0.508$ & $0.508$ \\
bio6     & $+0.291$ & -- & $0.250$ & $0.424$ \\
bio11    & $+0.219$ & -- & $0.396$ & $0.494$ \\
\bottomrule
\end{tabular}
\caption{Mitotype vs.\ ancestry/climate, 17 units. $\Omega$-null omitted where neither
variable is ancestry (see Erratum). $p_{\mathrm{BH}}$ = Benjamini-Hochberg corrected
across the 5 tests.}
\end{table}

\textbf{Finding.} Mitotype and genome-wide nuclear ancestry are nominally associated
($p\approx0.037$--$0.049$ uncorrected) but \textbf{this does not survive correction for the
5-comparison family} ($p_{\mathrm{BH}}=0.243$). No climate variable shows any evidence of a
mitotype association. The most defensible description is an \emph{association between
mitochondrial lineage and genome-wide nuclear ancestry, consistent with shared demographic
history} -- language such as ``founder-history signal'' or ``populations founded by an
\Faq-like queen'' asserts a specific causal mechanism (founding history) that this
correlational analysis cannot distinguish from alternatives such as asymmetric
introgression, retained population structure, or selection involving mitochondrial
ancestry.

\begin{figure}[H]\centering
  \includegraphics[width=\linewidth]{../Figures/moduleB_mitotype_vs_ancestry_climate_17units_CORRECTED.png}
  \caption{Mitotype vs.\ each variable, 17 independent lineage units.}
\end{figure}

\clearpage
\section*{\color{Accent}4. Does PC1's ancestry link survive controlling for mitotype?}

Since mitotype correlates with both ancestry and, more weakly, PC1, the PC1--ancestry
correlation could partly reflect the same demographic-history structure linking mitotype
to ancestry rather than an independent signal. Tested via partial correlation on the
17-unit dataset, with both nulls re-run on the partial statistic itself
($\Omega$-null draws correctly replacing ancestry; see Erratum).

\begin{table}[H]\centering\small
\begin{tabular}{lrr}
\toprule
 & $r$ (17 units) & Notes \\
\midrule
Raw cor(PC1, ancestry)                     & $-0.481$ & $\Omega$-null $p=0.044$; block-perm.\ $p=0.051$ \\
Partial cor(PC1, ancestry $\mid$ mitotype) & $-0.402$ & $\Omega$-null $p=0.097$; block-perm.\ $p=0.124$ \\
\bottomrule
\end{tabular}
\caption{PC1--ancestry correlation before/after controlling for mitotype, 17 units.}
\end{table}

\textbf{Finding.} Controlling for mitotype moderately attenuates the PC1--ancestry
correlation ($-0.481\rightarrow-0.402$), and the remaining association is not
distinguishable from the population-structure null under either test ($p=0.097$--$0.124$).
This is consistent with two readings -- a genuinely independent, if weaker, climate signal;
or partial confounding by mitotype-linked demographic history -- and \textbf{the data
cannot distinguish between them}. With only 17 independent lineage units, this sample
cannot separate the contributions of climate and maternal-lineage history with any
confidence; that is a limitation of the design, not a result in either direction.

\begin{figure}[H]\centering
  \includegraphics[width=0.85\linewidth]{../Figures/moduleB_PC1_ancestry_partial_mitotype_17units_CORRECTED.png}
  \caption{PC1 vs.\ ancestry, raw (left) and after removing the linear effect of mitotype
  from both variables (right), 17 units.}
\end{figure}

\clearpage
\section*{\color{Accent}5. Relationship to the BayPass locus-level scan}

The population-level tests above and the BayPass climate-association scan ask related but
\emph{different} questions. The population-level test asks whether a climate variable
correlates with \emph{one genome-wide ancestry summary} across populations. BayPass asks
whether \emph{particular loci or LD clusters} are unusually associated with that climate
variable, after accounting for population covariance ($\Omega$) and thousands of tests. A
population-level correlation being weak or non-significant does not, by itself, resolve
whether it is relevant to the locus-level scan -- the two statements below need to be held
together.

\textbf{The climate axes are moderately aligned with population structure, but that
alignment is not demonstrably stronger than expected by chance under that structure.} For
PC1, the ancestry correlation is moderately large after excluding \AA land ($r\approx-0.48$
on the corrected 17-unit dataset), but it is only suggestive under the $\Omega$-null
($p=0.044$), nominal under the corrected lineage permutation ($p\approx0.051$), and
non-significant after controlling for mitotype ($p=0.097$--$0.124$; Section~4). PC2 is
weaker after \AA land exclusion and clearly non-significant under both structure-aware
tests (Section~1).

\textbf{This is fully compatible with, and plausibly explains, both the original canonical
BayPass climate scan's null-calibrated result and the new Stage-1-direct scan's.} A
moderate alignment between climate and ancestry can cause loci that strongly differentiate
the underlying population groups to acquire large \emph{raw} Bayes factors, simply because
those loci track the same population split the climate variable happens to track. But
$\Omega$-structured null covariates -- random vectors sharing only the populations'
covariance structure, with no causal link to climate at all -- frequently produce
similarly-aligned, similarly-sized apparent outliers. Consequently the observed candidate
counts were not exceptional against that null, under \emph{either} scan:

\begin{table}[H]\centering\small
\begin{tabular}{lrrr}
\toprule
 & Stage-2 eMLG scan (32{,}854 clusters) & Stage-1-direct scan (18{,}361 units) & Expected by chance \\
\midrule
PC1 floor survivors & 0 & 3 & 1.84 (Stage-1) / 3.29 (Stage-2) \\
PC2 floor survivors & 4 & 3 & 1.84 (Stage-1) / 3.29 (Stage-2) \\
Set-level FDR & $\approx0.82$ (PC2) & $\approx0.61$ (both axes) & -- \\
\bottomrule
\end{tabular}
\caption{Floor-survivor candidate counts (BF(dB)$\geq$15 and beating all 10{,}000
$\Omega$-structured null covariates) under both BayPass scans. ``Expected by chance'' is
$N_{\mathrm{tested}}/(n_{\mathrm{sim}}+1)$.}
\end{table}

Both scans land at a high set-level FDR ($\geq0.6$) -- neither is a defensible discovery
set. The one notable shift is that PC1, which had \emph{zero} floor survivors under the
more heavily Stage-2-merged scan, has 3 under the less-diluted Stage-1-direct scan
(BF(dB) = 53.7, 38.7, 31.1) -- consistent with, though not confirming, the population-level
observation above that PC1 (not PC2) is the axis with the more suggestive ancestry link.
PC2's candidate count is essentially unchanged (4 vs.\ 3) and remains similarly
unremarkable relative to its own null expectation. The structured null is, in effect,
saying: \emph{these apparent climate associations are approximately what is expected when a
covariate happens to align with population history, regardless of how the markers were
reduced to test units.}

\textbf{Confounding does not have to be statistically significant to influence
locus-level results.} Statistical significance asks whether the observed alignment is
unusually strong given uncertainty and structure; confounding merely describes that climate
and ancestry point partly in the same direction. With only 17--19 independent populations, a
correlation of $0.4$--$0.5$ can be biologically noticeable, and large enough to shift which
loci look like outliers in a raw BayPass scan, while still being far too weak for a
population-level test with this little power to call it exceptional.

\textbf{Bottom line.} Both climate axes were partially aligned with genome-wide ancestry,
although these population-level correlations were not consistently distinguishable from
expectations under the population-covariance null. This moderate alignment nevertheless
provides a plausible source of the raw locus-level BayPass signal: loci that differentiate
the same population groups can appear climate-associated even without any locus-specific
environmental effect. Consistent with this interpretation, neither the number nor the
strength of cluster-level BayPass associations exceeded what was obtained under
$\Omega$-structured null covariates, under \emph{either} the Stage-2 or the Stage-1-direct
scan (set-level FDR $\geq0.6$ in both). The conclusion is \emph{not} that climate definitely
has no relationship with ancestry -- it is that the data cannot separate a climate effect
from population history, and the BayPass outliers provide no additional locus-specific
evidence beyond that structure. Even a formally significant genome-wide PC1--ancestry
correlation would not, on its own, rescue the BayPass candidates: it would show that overall
ancestry covaries with PC1, but evidence for particular climate-associated genomic regions
would still require their Bayes factors to exceed the structured-null distribution, which
they did not.

\clearpage
\section*{\color{Accent}6. Genome-wide architecture of the Stage-1-direct climate-association landscape}

Section~5 asked whether any \emph{individual} Stage-1 unit's climate Bayes factor is
exceptional against the $\Omega$-structured null (the floor-survivor test). A second,
complementary calibration asks a different, genome-wide question: across all 18{,}361
Stage-1-direct units, is the \emph{strength} of a unit's climate association --
its BF(dB) on PC1/PC2 -- itself predictable from genomic architecture (ancestry
differentiation, recombination rate, unidirectional ancestry sorting) any more or less
than an $\Omega$-structured pseudo-covariate with no causal link to climate would show by
chance? For each of PC1, PC2, and 10{,}000 null covariates, the genome-wide BF vector is
reduced to three statistics -- Spearman correlation with the per-unit Diagnostic Index
(DI), Spearman correlation with recombination rate, and a directional-sorting contrast
(mean BF percentile of unidirectionally-sorted vs.\ non-directional units, among
differentiated units only) -- and each observed statistic is compared to its own
10{,}000-value null distribution. The six resulting tests (2 axes $\times$ 3 architecture
measures) are Benjamini-Hochberg corrected as one family.

\begin{table}[H]\centering\small
\begin{tabular}{llrrrr}
\toprule
Test & Axis & Observed & Null median & Null 95\% & $p_{\mathrm{BH}}$ \\
\midrule
DI (Spearman $\rho$)            & PC1 & $+0.097$ & $+0.029$ & $[-0.018, 0.072]$ & $0.006$ \\
DI (Spearman $\rho$)            & PC2 & $-0.137$ & $+0.029$ & $[-0.018, 0.072]$ & $<0.001$ \\
Recombination (Spearman $\rho$) & PC1 & $-0.098$ & $+0.065$ & $[-0.013, 0.144]$ & $<0.001$ \\
Recombination (Spearman $\rho$) & PC2 & $+0.022$ & $+0.065$ & $[-0.013, 0.144]$ & $0.294$ \\
Sorting (differentiated only)   & PC1 & $+0.043$ & $+0.121$ & $[0.078, 0.165]$ & $0.001$ \\
Sorting (differentiated only)   & PC2 & $+0.019$ & $+0.121$ & $[0.078, 0.165]$ & $<0.001$ \\
\bottomrule
\end{tabular}
\caption{Genome-wide architecture calibration, Stage-1-direct universe (18{,}361 units).
``Null median''/``Null 95\%'' describe the same statistic computed on 10{,}000
$\Omega$-structured pseudo-covariates -- i.e.\ what a typical covariate sharing only the
populations' covariance structure, with no climate information, would produce.}
\end{table}

\textbf{Finding: five of the six tests survive FDR, but the directions do not tell a single
clean story.} DI is significantly correlated with BF strength on \emph{both} axes, in
\emph{opposite} directions ($+0.097$ on PC1, above the null range; $-0.137$ on PC2, below
it) -- consistent with PC1 and PC2 capturing different, partly opposing aspects of the same
underlying ancestry structure (Section~1 already found opposite-signed raw
ancestry--climate correlations for the two axes). Recombination shows a real effect only on
PC1: BF is \emph{more} concentrated in low-recombination regions than an $\Omega$-structured
covariate would produce by chance ($\rho=-0.098$, below the null), while PC2 is
indistinguishable from the null ($p_{\mathrm{BH}}=0.294$). Directional sorting is where the
result is least intuitive: on \emph{both} axes, real climate BF is \emph{less} elevated in
directionally-sorted (vs.\ non-directional) differentiated units than the null baseline
predicts ($0.043$ and $0.019$ vs.\ a null median of $0.121$) -- the \emph{opposite} of what a
naive ``sorted regions are more differentiated, so they mechanically inflate any
covariate's BF'' argument would predict.

This genome-wide result should not be over-read as either confirming or contradicting
Section~5's conclusion. It does not identify individual climate-adaptive loci -- that
remains the floor-survivor test's job, and it still returns a high set-level FDR ($\geq0.6$)
under both scans. What it adds is a more granular picture: the Stage-1-direct climate-BF
landscape is \emph{not} architecture-agnostic (5 of 6 tests depart from the structured-null
baseline), but the departures point in different, and for sorting an unexpected, direction
on the two climate axes -- so \emph{which} architectural feature aligns with which axis, and
whether that alignment reflects the same ancestry confound discussed in Sections~1--4 or a
distinct mechanism, is not resolved here and would need its own follow-up (e.g.\ relating
each axis's DI/recombination/sorting correlation directly to that axis's own
ancestry-confound coefficient from \texttt{moduleB\_ancestry\_confound.R}).

\clearpage
\section*{\color{Accent}Limitations and recommended follow-up}\label{sec:limitations}
This analysis should be treated as \textbf{supplementary and explicitly exploratory}, not a
confirmed result, for several reasons:
\begin{itemize}
  \item \textbf{Only PC1 received the full mitotype-controlled (partial) analysis.} PC2,
    bio6, and bio11 are examined only through bivariate correlations. The conclusion that no
    climate variable shows an ancestry- and mitotype-independent signal is broader than what
    has actually been tested for those three variables; either the same partial model should
    be run for all four, or the conclusion should be narrowed to PC1 specifically.
  \item \textbf{Ancestry was computed once, from the frozen DI$>-25$ marker panel.} It has
    not been recomputed from an independent marker set or from LD-reduced diagnostic units,
    so its sensitivity to the specific marker panel used is untested.
  \item \textbf{A single coherent structure-aware model is preferable to a sequence of
    bivariate/partial tests.} A model of the form
    $\mathrm{ancestry} \sim \mathrm{climate} + \mathrm{mitotype}$ with an $\Omega$-based
    residual covariance (or an equivalent parametric bootstrap) would estimate both
    predictors' contributions jointly, with correctly propagated uncertainty, rather than
    inferring their relative roles from separate bivariate/partial tests as done here.
  \item With 17 independent lineage units and 4-5 tests per family, this analysis is
    underpowered to detect anything but a fairly strong effect, and $\Omega$-null and
    block-permutation $p$-values disagree by up to $\sim$0.03--0.05 in places (e.g.\ PC1:
    $0.044$ vs.\ $0.051$) -- a reminder that ``significant'' here would depend on which
    defensible null is treated as primary, not a settled question either null can resolve
    alone.
  \item \textbf{Section~6's genome-wide architecture calibration has mixed-direction results
    that were not chased down.} DI, recombination, and sorting depart from the
    $\Omega$-structured null on different axes and, for sorting, in the opposite direction
    from a naive confounding argument. Relating each axis's departure directly to that
    axis's own ancestry-confound coefficient (rather than to the coarser population-level
    ancestry summary used in Sections~1--4) would be needed before drawing a mechanistic
    conclusion from it.
\end{itemize}

\section*{\color{Accent}Summary}
Across the 19 \AA land-excluded populations (17 independent lineage units after correcting
for shared origins), genome-wide ancestry showed a moderate negative correlation with
climate PC1 ($r=-0.48$). This relationship was nominally suggestive under the
shared-lineage permutation ($p=0.051$) but did not clear the $\Omega$-structured null
($p=0.044$ is itself only marginal, and the two nulls disagree), and it did not survive
correction for testing four climate variables ($p_{\mathrm{BH}}=0.16$). The association was
further attenuated and statistically indistinguishable from the population-structure null
after controlling for mitotype ($r=-0.40$, $p=0.10$--$0.12$). Mitotype was itself nominally
associated with nuclear ancestry ($p\approx0.04$) but not with any measured climate
variable, and that association also did not survive multiple-testing correction. These
patterns suggest contributions from population history and possibly climate, but the
limited number of independent populations does not permit their effects to be separated
confidently. Critically, this moderate, inconclusive population-level alignment is
sufficient to explain the raw candidate counts in the original BayPass climate scan without
invoking any locus-specific climate effect (Section~5): no PC1 cluster and only a
null-expected number of PC2 clusters exceeded the $\Omega$-structured null. The BayPass
scan therefore provides no locus-specific evidence for climate adaptation beyond what
population structure already predicts. A separate, genome-wide calibration (Section~6)
shows the Stage-1-direct climate-BF landscape is not architecture-agnostic -- five of six
tests relating BF strength to Diagnostic Index, recombination, and directional sorting
depart from the structured-null baseline -- but the directions differ across the two
climate axes and, for sorting, run counter to a simple confounding story, so this result
sharpens rather than resolves the population-level picture above. This document should be
read as a supplementary, exploratory analysis, not a confirmed finding for the main text.

\end{document}
